% TOP_PROBLEM: {"id": 30007095, "title": "Integral and fractional expectation thresholds within a constant factor", "category_id": 2, "category": {"id": 2, "name": "combinatorics", "display_name": "Combinatorics", "description": "Counting problems, graph theory, discrete structures.", "slug": "combinatorics", "order_index": 2, "created_at": "2026-07-31T15:26:25.671Z"}, "status": "open"}
% FIRST_POSTED: 2026-10-06
% ATTEMPT_STATUS: solved
% OPENAI_CLAIM: {"schema_version":1,"source_repo":"https://github.com/openai/math","source_commit":"adc7f1241b42e322a6451854ab7e4b4c146bf78a","release_date":"2026-10-06","match":"stronger","resolution":"proved","independently_reviewed":false,"second_pass":"agreed","families":[{"family":"175","title":"Talagrand’s expectation thresholds, discrete convexity, and graph decompositions","subject":"Combinatorics","manuscripts":[{"title":"Integral and fractional expectation thresholds are equivalent","pdf_path":"preprints/Integral-and-fractional-expectation-thresholds-are-equivalent-September-23-2026/paper.pdf","readme_path":"preprints/Integral-and-fractional-expectation-thresholds-are-equivalent-September-23-2026/README.md","date":"2026-09-23","theorem_ref":"thm:main (source/sections/introduction.tex:52)","statement_sources":[{"path":"preprints/Integral-and-fractional-expectation-thresholds-are-equivalent-September-23-2026/build/source/sections/introduction.tex","line":52,"label":"thm:main"}]}],"lean":{"doc_path":"lean/docs/175.md","comparators":["lean/ComparatorChallenges/TalagrandDiscreteConvexity.json","lean/ComparatorChallenges/TalagrandDiscreteConvexity.lean","lean/ComparatorChallenges/TalagrandExpectationThreshold.json","lean/ComparatorChallenges/TalagrandExpectationThreshold.lean"],"declarations":["OAI.TalagrandDiscreteConvexity.talagrand_discrete_convexity","OAI.TalagrandThreshold.talagrand_expectation_threshold_equivalence"],"covers_main_theorem":true},"reasoning_trace":null}]}
\documentclass[11pt]{article}
\usepackage[margin=25mm]{geometry}
\usepackage{fontspec}
\setmainfont{Latin Modern Roman}
\usepackage{amsmath,amssymb,mathtools}
\usepackage{unicode-math}
\setmathfont{Latin Modern Math}
\usepackage{enumitem}
\usepackage{xurl,hyperref,bookmark}
\hypersetup{colorlinks=true,urlcolor=blue}
\setcounter{secnumdepth}{0}
\setlength{\parindent}{0pt}
\setlength{\parskip}{5pt}
\setlength{\emergencystretch}{3em}
\begin{document}
\section*{Integral and fractional expectation thresholds within a constant factor}

\subsection{OpenAI's claim}
\textit{Integral versus fractional thresholds}: The claimed comparison bounds the fractional expectation threshold by a universal multiple of its integral version. Both versions are defined with the same containment-cover budget.

\subsection{Mathematical statement of the claimed result}
\paragraph{Mathematical claim 1.} For nonempty proper increasing $\mathcal F\subset2^X$, with finite nonempty $X$, let $q(\mathcal F)$ be the greatest $p\in[0,1]$ admitting a cover $\mathcal G$ such that every $H\in\mathcal F$ contains some $S\in\mathcal G$ and $\sum_{S\in\mathcal G}p^{|S|}\le1/2$. Define $q_f$ by allowing weights $g:2^X\to[0,1]$, requiring $\sum_{S\subseteq H}g(S)\ge1$ for each $H\in\mathcal F$ and $\sum_Sg(S)p^{|S|}\le1/2$. The exact claim is \[q(\mathcal F)\le q_f(\mathcal F)\le25\cdot512^4q(\mathcal F).\] The same budget $1/2$ is used on both sides, and the factor is independent of the ground-set size and minimal-member sizes.

\subsection{Outline of the argument}
\textit{Integral versus fractional thresholds}: A multiscale random selector finds one member of the family that captures enough mass at every scale at once. Normalizing fractional cover weights on vertices gives a bounded average color for that member. Independent exponential color variables then force incompatible upper and lower expectations when the integral threshold is too small relative to the fractional threshold.

\subsection{Scope relative to this problem}
The theorem covers every nonempty proper increasing family on a finite ground set with both cover budgets equal to one half. Its explicit universal factor $25\cdot512^4$ answers the entire constant-factor question. The empty-ground-set case has no such family.

\subsection{Formal verification}
Family 175 describes Lean coverage for the integral/fractional threshold comparison and the discrete-convexity statement; graph decomposition is a later separate consequence.

\textit{Integral versus fractional thresholds}: The manuscript links \texttt{TalagrandDiscreteConvexity.lean} (declaration \texttt{OAI.TalagrandDiscreteConvexity.talagrand\_discrete\_convexity}), \texttt{TalagrandExpectationThreshold.lean} (declaration \texttt{OAI.TalagrandThreshold.talagrand\_expectation\_threshold\_equivalence}).

The documented selection includes this manuscript’s main assertion; that assertion retains the restrictions stated above.

This site has not run the Lean development. These links describe OpenAI’s selected formal statements and should be read with their stated scope.

\subsection{Status caveat}
OpenAI reports that an internal model produced the manuscripts. Their arguments have not been independently verified as proofs by this site. Any portion outside the documented formal scope may contain mistakes, and the catalogue label reports the announced mathematical claim.

\subsection{Sources}
\begin{itemize}
\item OpenAI, \emph{Integral and fractional expectation thresholds are equivalent}, 2026-09-23; citation key \texttt{OAI:Integral-and-fractional-expectation-thresholds-are-equivalent-September-23-2026}. \href{https://github.com/openai/math/blob/adc7f1241b42e322a6451854ab7e4b4c146bf78a/preprints/Integral-and-fractional-expectation-thresholds-are-equivalent-September-23-2026/paper.pdf}{PDF (GitHub)}; \href{https://github.com/openai/math/raw/adc7f1241b42e322a6451854ab7e4b4c146bf78a/preprints/Integral-and-fractional-expectation-thresholds-are-equivalent-September-23-2026/paper.pdf}{PDF (direct)}; \href{https://github.com/openai/math/blob/adc7f1241b42e322a6451854ab7e4b4c146bf78a/preprints/Integral-and-fractional-expectation-thresholds-are-equivalent-September-23-2026/README.md}{Manuscript page and citation}; \href{https://github.com/openai/math/blob/main/preprints/Integral-and-fractional-expectation-thresholds-are-equivalent-September-23-2026/paper.pdf}{Latest version}; Statement source: \href{https://github.com/openai/math/blob/adc7f1241b42e322a6451854ab7e4b4c146bf78a/preprints/Integral-and-fractional-expectation-thresholds-are-equivalent-September-23-2026/build/source/sections/introduction.tex\#L52}{thm:main}.
\item Family 175: \href{https://github.com/openai/math/blob/adc7f1241b42e322a6451854ab7e4b4c146bf78a/lean/docs/175.md}{Lean scope}.
\item \href{https://github.com/openai/math/blob/adc7f1241b42e322a6451854ab7e4b4c146bf78a/lean/ComparatorChallenges/TalagrandDiscreteConvexity.json}{Comparator TalagrandDiscreteConvexity.json}.
\item \href{https://github.com/openai/math/blob/adc7f1241b42e322a6451854ab7e4b4c146bf78a/lean/ComparatorChallenges/TalagrandDiscreteConvexity.lean}{Comparator TalagrandDiscreteConvexity.lean}.
\item \href{https://github.com/openai/math/blob/adc7f1241b42e322a6451854ab7e4b4c146bf78a/lean/ComparatorChallenges/TalagrandExpectationThreshold.json}{Comparator TalagrandExpectationThreshold.json}.
\item \href{https://github.com/openai/math/blob/adc7f1241b42e322a6451854ab7e4b4c146bf78a/lean/ComparatorChallenges/TalagrandExpectationThreshold.lean}{Comparator TalagrandExpectationThreshold.lean}.
\item \href{https://github.com/openai/math/blob/adc7f1241b42e322a6451854ab7e4b4c146bf78a/CONTENTS.md}{OpenAI catalogue entry}.
\end{itemize}
\end{document}
